Los canales de comunicación del proyecto definen los mecanismos mediante los cuales los miembros del equipo y los diferentes interesados intercambian información durante el ciclo de vida del desarrollo del software. Estos canales se establecen en coherencia con el Protocolo de Entrega definido en la sección 10.2, el cual combina comunicación interactiva, comunicación \textit{push} y comunicación \textit{pull} para la distribución y validación de los entregables del proyecto.

\textbf{Equipo de Dirección del Proyecto}
\begin{itemize}
\item Project Manager
\item Líder Técnico
\item Patrocinador del proyecto
\item Comité Directivo
\end{itemize}

Este equipo es responsable de la toma de decisiones estratégicas, la validación de entregables clave y el escalamiento de incidentes que puedan afectar el avance del proyecto.

\textbf{Equipo de Desarrollo}
\begin{itemize}
\item Desarrollador Backend
\item Desarrollador Frontend
\item Analista QA
\item Administrador de Base de Datos / DevOps
\end{itemize}

Este equipo mantiene la comunicación técnica relacionada con el desarrollo del sistema, la implementación de módulos, el control del repositorio y la resolución de incidencias técnicas.

\textbf{Equipo Operativo del Cliente}
\begin{itemize}
\item Operadores de ventanilla
\item Supervisores de tránsito
\item Área financiera / tesorería
\item Área de TI institucional
\end{itemize}

Este grupo participa principalmente en la validación funcional de los incrementos de software durante las revisiones de sprint y en la retroalimentación sobre el funcionamiento del sistema.

\textbf{Entidades Externas}
\begin{itemize}
\item Ministerio de Transporte
\item RUNT
\item Proveedor de impresión
\item Entidades regulatorias
\end{itemize}

Estas entidades participan en procesos de validación técnica, integración del sistema y cumplimiento de normativas.

\subsubsection{Reuniones Tecnicas periódicas}

\begin{itemize}

\item \textbf{Reunión de planificación del Sprint} \\
Frecuencia: inicio de cada sprint. \\
Participantes: Project Manager y equipo de desarrollo. \\
Objetivo: definir las tareas y entregables del sprint.

\item \textbf{Reunión de seguimiento} \\
Frecuencia: semanal. \\
Participantes: equipo de desarrollo y Project Manager. \\
Objetivo: revisar el progreso del proyecto, identificar bloqueos y coordinar actividades.

\item \textbf{Reunión de demostración (Sprint Review)} \\
Frecuencia: al finalizar cada sprint. \\
Participantes: equipo de desarrollo, patrocinador y usuarios clave. \\
Objetivo: presentar el incremento de software desarrollado y recibir retroalimentación.

\item \textbf{Reuniones técnicas de integración} \\
Frecuencia: según necesidad. \\
Participantes: equipo técnico e integradores externos. \\
Objetivo: validar integraciones tecnológicas y resolver problemas técnicos.

\item \textbf{Reunión de cierre del proyecto} \\
Frecuencia: al finalizar el proyecto. \\
Participantes: todos los stakeholders principales. \\
Objetivo: presentar el sistema final y formalizar la entrega del proyecto.

\end{itemize}

\subsubsection{2.3 Medios físicos}

Aunque el proyecto utiliza principalmente medios digitales, se contemplan algunos medios físicos para formalizar documentación institucional y presentaciones oficiales.

\begin{itemize}
\item Actas impresas de reuniones formales cuando sea requerido.
\item Documentación contractual o administrativa firmada físicamente.
\item Informes técnicos impresos utilizados en presentaciones institucionales.
\item Demostraciones presenciales del sistema para validaciones operativas en las entidades de tránsito.
\end{itemize}

\subsubsection{2.4 Correos electrónicos, llamadas y videollamadas}

Estos canales constituyen los principales mecanismos de comunicación operativa del proyecto y permiten mantener la coordinación entre los equipos durante el desarrollo del software.

\textbf{Correo electrónico}
De acuerdo con el protocolo del proyecto, todo correo enviado requiere confirmación de recepción por parte del destinatario en un plazo máximo de 24 horas.

Se utilizará principalmente para:

\begin{itemize}
\item Envío de actas de reuniones.
\item Entrega de documentos formales del proyecto.
\item Notificación de disponibilidad de nuevos incrementos de software.
\item Recepción de observaciones de los interesados.
\end{itemize}



\textbf{Videollamadas}

Se utilizan para reuniones interactivas del proyecto, especialmente para:

\begin{itemize}
\item demostraciones de módulos desarrollados,
\item revisiones de sprint,
\item reuniones técnicas con integradores externos,
\item reuniones del comité directivo.
\end{itemize}

\textbf{Llamadas telefónicas}

Las llamadas telefónicas se utilizan principalmente para la coordinación inmediata entre miembros del equipo y para la resolución de incidentes urgentes que puedan afectar el desarrollo del proyecto.

\textbf{Repositorios en la nube}

Para facilitar la colaboración entre los miembros del equipo y garantizar la disponibilidad permanente de la información del proyecto, se utilizarán herramientas de almacenamiento y gestión de repositorios en la nube. Estas plataformas permitirán centralizar los artefactos del proyecto y asegurar que todos los interesados puedan acceder a la versión más actualizada de la documentación y del software.

Esto facilita la coordinación entre los diferentes roles del proyecto y evita la duplicación o pérdida de información.

Dentro de este entorno se almacenarán y gestionarán los siguientes elementos del proyecto:

\begin{itemize}
\item Código fuente del sistema y control de versiones del software.
\item Documentación técnica del proyecto (arquitectura, diagramas y especificaciones).
\item Artefactos de diseño y modelos del sistema.
\item Manuales técnicos y material de capacitación.
\item Registros de pruebas, reportes de errores y documentación de despliegue.
\end{itemize}

El uso de estas herramientas permite que el equipo de desarrollo trabaje de manera colaborativa, manteniendo trazabilidad de los cambios realizados en el proyecto y facilitando la integración continua de los diferentes módulos del sistema.